\chapterbackmatter{Solutions to Supplementary Exercises}

\begin{enumerate}
\SupplementarySolution{1}{Invariant Mass-Shell Measure}
For a proper Lorentz transformation,
\(\lvert\det\Lambda\rvert=1\), so \(d^4p\) is invariant; \(p^2\) is a
scalar, so is \(\delta(p^2+m^2)\).  On the positive-energy shell,
\[
 \delta(p^2+m^2)
 =\delta\!\left(-(p^0)^2+\mathbf p^2+m^2\right)
 =\frac{\delta(p^0-E_{\mathbf p})+\delta(p^0+E_{\mathbf p})}
 {2E_{\mathbf p}}.
\]
Multiplication by \(\theta(p^0)\) selects the first root, and hence
\[
 \int d^4p\,\delta(p^2+m^2)\theta(p^0)F(p)
 =\int\frac{d^3\mathbf p}{2E_{\mathbf p}}F(E_{\mathbf p},\mathbf p).
\]
A proper orthochronous transformation maps every future-directed causal
vector to another such vector.  It therefore preserves the sign of \(p^0\).
Time reversal interchanges the two energy sheets, so the positive-sheet
measure alone is not invariant under that disconnected component.

\SupplementarySolution{2}{The Invariant Delta Kernel}
The invariant on-shell measure is \(d\mu(p)=d^3\mathbf p/E_{\mathbf p}\),
apart from an irrelevant constant.  Its delta kernel \(\Delta(p,q)\) is
defined by
\[
 f(p)=\int d\mu(q)\,\Delta(p,q)f(q).
\]
Comparison with the ordinary three-dimensional delta function gives
\[
 \Delta(p,q)=E_{\mathbf q}\delta^3(\mathbf q-\mathbf p)
            =E_{\mathbf p}\delta^3(\mathbf p-\mathbf q).
\]
Because both \(d\mu\) and the defining identity are invariant, \(\Delta\)
is invariant.

The common covariant convention is
\[
 \braket{p',\sigma'}{p,\sigma}_{\rm cov}
 =2E_{\mathbf p}(2\pi)^3
 \delta_{\sigma'\sigma}\delta^3(\mathbf p-\mathbf p').
\]
The convention of Eq.~\eqref{eq:2.5.19} omits the factor
\(2E_{\mathbf p}(2\pi)^3\).  The compensating factor
\(\sqrt{(\Lambda p)^0/p^0}\) then appears in the state transformation law.

\SupplementarySolution{3}{A Lorentz-Normalized Wave Packet}
Use
\[
 \braket{\phi}{\psi}
 =\int\frac{d^3\mathbf p}{2E_{\mathbf p}}\,
   \phi(\mathbf p)^*\psi(\mathbf p).
\]
For a scalar representation,
\[
 (U(\Lambda)\psi)(\mathbf p)
 =\psi(\boldsymbol{\Lambda^{-1}p}),
\]
where the energy component of \(p\) is fixed by the positive mass shell.
Then, with \(q=\Lambda^{-1}p\),
\[
 \|U(\Lambda)\psi\|^2
 =\int\frac{d^3\mathbf p}{2E_{\mathbf p}}
   |\psi(\boldsymbol{\Lambda^{-1}p})|^2
 =\int\frac{d^3\mathbf q}{2E_{\mathbf q}}|\psi(\mathbf q)|^2.
\]
The last equality is precisely the invariance established in S.2.1.
Thus the scalar pullback is unitary without an additional Jacobian factor.

\SupplementarySolution{4}{Boost Jacobian on Shell}
For the passive boost
\[
 E'=\gamma(E-up_z),\qquad
 p'_z=\gamma(p_z-uE),\qquad p'_x=p_x,\quad p'_y=p_y.
\]
Since \(\partial E/\partial p_z=p_z/E\),
\[
 \frac{\partial p'_z}{\partial p_z}
 =\gamma\left(1-u\frac{p_z}{E}\right)
 =\frac{\gamma(E-up_z)}{E}=\frac{E'}{E}.
\]
The transverse coordinates are unchanged and their mixed derivatives do
not alter the triangular determinant, so
\[
 \det\frac{\partial\mathbf p'}{\partial\mathbf p}
 =\frac{E'}{E}.
\]
Consequently \(d^3\mathbf p'=(E'/E)d^3\mathbf p\), which immediately yields
\(d^3\mathbf p'/E'=d^3\mathbf p/E\).

\SupplementarySolution{5}{Poincar\'e Algebra from Differential Operators}
The elementary relations
\[
 [\partial^\mu,x^\nu]=\eta^{\mu\nu},\qquad
 [\partial^\mu,\partial^\nu]=[x^\mu,x^\nu]=0
\]
give \([P^\mu,P^\nu]=0\).  Applying the product rule once gives
\[
 [J^{\mu\nu},P^\rho]
 =i\bigl(\eta^{\mu\rho}P^\nu-\eta^{\nu\rho}P^\mu\bigr),
\]
which is equivalent to Eq.~\eqref{eq:2.4.13}.  Applying it twice gives
\[
\begin{split}
 [J^{\mu\nu},J^{\rho\sigma}]
 =i(&\eta^{\mu\rho}J^{\nu\sigma}
 -\eta^{\nu\rho}J^{\mu\sigma}\\
 &+\eta^{\nu\sigma}J^{\mu\rho}
 -\eta^{\mu\sigma}J^{\nu\rho}),
\end{split}
\]
equivalent to Eq.~\eqref{eq:2.4.12}.  No field equation was used: these
relations express the action of infinitesimal spacetime isometries on
scalar functions.

\SupplementarySolution{6}{Rotations and Boosts as a Closed Algebra}
Substituting
\(J_i=\frac12\epsilon_{ijk}J^{jk}\), \(K_i=J^{i0}\), and \(H=P^0\)
into the covariant algebra yields
\[
\begin{gathered}
 [J_i,J_j]=i\epsilon_{ijk}J_k,\qquad
 [J_i,K_j]=i\epsilon_{ijk}K_k,\qquad
 [K_i,K_j]=-i\epsilon_{ijk}J_k,\\
 [J_i,P_j]=i\epsilon_{ijk}P_k,\qquad
 [J_i,H]=0,\qquad
 [K_i,P_j]=i\delta_{ij}H,\qquad
 [K_i,H]=iP_i,\\
 [P_i,P_j]=[P_i,H]=0.
\end{gathered}
\]
The minus sign in \([K_i,K_j]\) is the algebraic signature of the
noncompact Lorentz boosts.  The relations involving \(H\) also show why a
boost does not preserve an energy eigenspace.

\SupplementarySolution{7}{Scalar Generators in Momentum Space}
Multiplication by real \(p^\mu\) is Hermitian.  The rotation vector fields
preserve both \(\mathbf p^2\) and \(d^3\mathbf p\), so integration by parts
makes \(\mathbf J\) Hermitian.  For a boost,
\[
 \frac{d^3\mathbf p}{2E}\,E\partial_{p_i}
 =\frac12d^3\mathbf p\,\partial_{p_i};
\]
the boundary term vanishes for smooth rapidly decreasing functions, and
\(iE\partial_{p_i}\) is Hermitian.

The useful direct checks are
\[
\begin{aligned}
 [K_i,P_j]\psi&=iE\partial_i(p_j\psi)-p_j iE\partial_i\psi
   =i\delta_{ij}E\psi,\\
 [K_i,H]\psi&=iE(\partial_iE)\psi=ip_i\psi,\\
 [K_i,K_j]\psi
 &=-\bigl(p_i\partial_j-p_j\partial_i\bigr)\psi
 =-i\epsilon_{ijk}J_k\psi.
\end{aligned}
\]
The remaining brackets follow from the usual orbital angular-momentum
algebra.  Hence these operators realize the spin-zero one-particle
Poincar\'e representation.

\SupplementarySolution{8}{Finite Translations of Momentum Eigenstates}
From Eq.~\eqref{eq:2.4.3},
\[
 U(\mathbf1,a)=\exp(-ia_\mu P^\mu),
\]
and therefore
\[
 U(\mathbf1,a)\ket{p,\sigma}
 =e^{-ia_\mu p^\mu}\ket{p,\sigma}.
\]
Using \(U(\Lambda)P^\mu U(\Lambda)^{-1}
=\Lambda_\nu{}^\mu P^\nu\),
\[
\begin{split}
 U(\Lambda)U(\mathbf1,a)U(\Lambda)^{-1}
 &=\exp\!\left[-ia_\mu\Lambda_\nu{}^\mu P^\nu\right]\\
 &=\exp[-i(\Lambda a)_\nu P^\nu]
 =U(\mathbf1,\Lambda a).
\end{split}
\]
This is the operator form of the semidirect-product multiplication law.

\SupplementarySolution{9}{Pauli--Lubanski Orthogonality}
Antisymmetry alone gives
\[
 \mathcal W_\mu P^\mu
 =\frac12\epsilon_{\mu\nu\rho\sigma}
 J^{\nu\rho}P^\sigma P^\mu=0,
\]
because the momentum operators commute, making
\(P^\sigma P^\mu\) symmetric.  Next,
\[
 [\mathcal W_\mu,P^\lambda]
 =\frac12\epsilon_{\mu\nu\rho\sigma}
 [J^{\nu\rho},P^\lambda]P^\sigma.
\]
Inserting the Poincar\'e commutator produces two terms, each containing a
symmetric product of momenta contracted with an antisymmetric pair of
indices.  They cancel, so
\([\mathcal W_\mu,P^\lambda]=0\).

\SupplementarySolution{10}{Spin from the Rest-Frame Casimir}
At \(p^\mu=(m,\mathbf0)\), the antisymmetry of
\(\epsilon_{\mu\nu\rho\sigma}\) gives
\[
 \mathcal W^0=0,\qquad
 \boldsymbol{\mathcal W}=\pm m\mathbf J,
\]
where the overall sign depends on the placement of indices in the
Levi--Civita convention and is irrelevant to the square.  Thus, on a
spin-\(j\) rest state,
\[
 \mathcal W^2=m^2\mathbf J^2=m^2j(j+1).
\]
Since S.2.11 shows that \(\mathcal W^2\) is a Lorentz scalar and a
Poincar\'e Casimir, its eigenvalue is the same in every inertial frame and
throughout an irreducible massive multiplet.

\SupplementarySolution{11}{The Pauli--Lubanski Pseudovector}
Under a translation,
\[
 J^{\nu\rho}\mapsto J^{\nu\rho}-a^\nu P^\rho+a^\rho P^\nu.
\]
The induced change in \(\mathcal W_\mu\) vanishes because it is
proportional to \(\epsilon_{\mu\nu\rho\sigma}P^\rho P^\sigma\).
For a Lorentz transformation, substituting the tensor law for \(J\), the
vector law for \(P\), and the determinant identity for the Levi--Civita
tensor gives
\[
 \mathcal W'_\mu=(\det\Lambda)\Lambda_\mu{}^\nu\mathcal W_\nu.
\]
It is a vector for proper transformations and changes the additional sign
appropriate to a pseudovector for improper ones.  In either case
\(\mathcal W'^2=\mathcal W^2\).  Infinitesimally this says
\([\mathcal W^2,P^\mu]=[\mathcal W^2,J^{\mu\nu}]=0\).

\SupplementarySolution{12}{What the Two Casimirs Do and Do Not Distinguish}
For a massive positive-energy representation,
\[
 P^2=-m^2,\qquad
 \mathcal W^2=m^2j(j+1).
\]
Thus \(j=0,\frac12,1\) give respectively
\[
 \mathcal W^2=0,\qquad
 \mathcal W^2=\frac34m^2,\qquad
 \mathcal W^2=2m^2.
\]
For a finite-helicity massless representation the null translations of
the little group vanish.  The remaining relation is
\(\mathcal W^\mu=\lambda P^\mu\), and hence
\[
 P^2=0,\qquad \mathcal W^2=\lambda^2P^2=0
\]
for every \(\lambda\).  The missing datum is the one-dimensional
representation of the rotational part of the little group---the helicity
\(\lambda\), including its sign.

\SupplementarySolution{13}{The Massive Little Group}
If \(\Lambda k=k\) for \(k=(m,\mathbf0)\), then
\(\Lambda^\mu{}_0=\delta^\mu{}_0\).  The Lorentz condition forces
\(\Lambda^0{}_i=0\) and leaves a spatial matrix \(R\) satisfying
\(R^{\mathsf T}R=\mathbf1\).  Properness gives \(\det R=1\), so
\[
 \Lambda=\begin{pmatrix}1&0\\0&R\end{pmatrix},
 \qquad R\in SO(3).
\]
The unitary irreducible representations of its double cover \(SU(2)\) are
labeled by \(j=0,\frac12,1,\ldots\) and have dimension \(2j+1\).  Inducing
each such little-group representation over the positive mass shell gives
the states \(\ket{p,\sigma}\), \(\sigma=-j,\ldots,j\).

\SupplementarySolution{14}{The Null Little-Group Algebra}
Acting with the vector-representation matrices on
\(k=(\kappa,0,0,\kappa)\) shows
\((J_2+K_1)k=(-J_1+K_2)k=J_3k=0\).  The Lorentz algebra gives
\[
 [J_3,A]=iB,\qquad
 [J_3,B]=-iA,\qquad
 [A,B]=0.
\]
Thus \(A,B\) are two commuting translations rotated into one another by
\(J_3\).  This is the Lie algebra
\(\mathfrak{iso}(2)=\mathbb R^2\rtimes\mathfrak{so}(2)\), the algebra of
the two-dimensional Euclidean group.  It is the connected little group of
a null momentum in \(3+1\) dimensions.

\SupplementarySolution{15}{Why Finite Helicity Sets Null Translations to Zero}
Choose a simultaneous eigenstate
\(\ket{a,b}\).  A rotation through \(\theta\) about the momentum sends its
eigenvalue pair to
\[
 (a_\theta,b_\theta)
 =(a\cos\theta+b\sin\theta,-a\sin\theta+b\cos\theta).
\]
If \(\rho^2=a^2+b^2>0\), these pairs sweep a full circle.  Distinct
eigenvalues of commuting Hermitian operators label orthogonal generalized
eigenstates, so the representation contains a continuous infinity of
polarizations.  A finite-dimensional polarization space is possible only
for \(\rho=0\), which means \(A=B=0\).  The remaining \(SO(2)\) character
is \(e^{i\lambda\theta}\) and supplies the helicity.

\SupplementarySolution{16}{The Wigner Phase of a Helicity State}
Let
\[
 W(\Lambda,p)=L^{-1}(\Lambda p)\Lambda L(p).
\]
When null translations act trivially, only the rotation angle
\(\theta(\Lambda,p)\) of this little-group element remains.  With the
\(\delta^3\)-normalized states of Chapter 2,
\[
 U(\Lambda)\ket{p,\lambda}
 =\sqrt{\frac{(\Lambda p)^0}{p^0}}\,
 e^{i\lambda\theta(\Lambda,p)}
 \ket{\Lambda p,\lambda}.
\]
For a spatial rotation by \(\theta\) about \(\mathbf p\), the standard
boost may be chosen so that \(W\) is the same rotation.  Momentum is
unchanged and the state acquires \(e^{i\lambda\theta}\), confirming that
\(\lambda\) is the angular momentum along the direction of motion.

\SupplementarySolution{17}{Parity, Time Reversal, and Helicity}
Parity sends
\[
 \mathbf p\mapsto-\mathbf p,\qquad
 \mathbf J\mapsto\mathbf J,
\]
because angular momentum is axial.  Therefore
\(P:\lambda\mapsto-\lambda\).  A parity-invariant spectrum containing a
nonzero-helicity particle must also contain an equal-mass
opposite-helicity partner.

Time reversal is antiunitary and sends both
\[
 \mathbf p\mapsto-\mathbf p,\qquad
 \mathbf J\mapsto-\mathbf J.
\]
Hence \(T:\lambda\mapsto\lambda\).  Time-reversal invariance can act within
a single helicity representation and does not by itself require the
opposite helicity.  These conclusions concern helicity; additional
internal quantum numbers may of course also be transformed.

\SupplementarySolution{18}{Normalization from a Standard Boost}
Unitarity first gives
\[
 \braket{p',\sigma'}{p,\sigma}
 =|N(p)|^2\delta_{\sigma'\sigma}
 \frac{p^0}{k^0}\delta^3(\mathbf p'-\mathbf p).
\]
Thus \(\delta^3\)-normalization requires
\[
 N(p)=\sqrt{\frac{k^0}{p^0}},
\]
up to a momentum-independent phase.  For a massive standard momentum
\(k^0=m\); for a massless standard momentum \(k^0=\kappa\), whose arbitrary
fixed value cancels from physical transformation laws.

Applying \(U(\Lambda)\), inserting
\(\mathbf1=U(L(\Lambda p))U(L(\Lambda p))^{-1}\), and using the ratio of
the two normalization factors gives
\[
 \frac{N(p)}{N(\Lambda p)}
 =\sqrt{\frac{(\Lambda p)^0}{p^0}}.
\]
This is the square-root multiplying the little-group matrix in
Eq.~\eqref{eq:2.5.23}.

\SupplementarySolution{19}{A Wigner Rotation from Perpendicular Boosts}
Represent the two boosts by their rapidity matrices and factor
\[
 B_z(\eta)B_y(\xi)=L(p_{\rm final})R_x(\Omega),
\]
where \(L(p_{\rm final})\) is the unique rotation-free boost taking the
rest momentum to the final momentum.  Comparing the off-diagonal entries
(equivalently, using the polar decomposition of the corresponding
\(SL(2,\mathbb C)\) matrices) gives
\[
 \tan\frac{\Omega}{2}
 =\frac{\sinh(\eta/2)\sinh(\xi/2)}
 {\cosh(\eta/2)\cosh(\xi/2)}
 =\tanh\frac{\eta}{2}\tanh\frac{\xi}{2}.
\]
The axis is parallel to the oriented cross product of the second and first
boost directions, here the \(x\)-axis up to convention.  Replacing an
active boost of states by the passive boost of coordinates inverts one
boost and therefore reverses \(\Omega\), but leaves every mixing
probability unchanged.

\SupplementarySolution{20}{Spin-One Mixing under the Wigner Rotation}
The absolute squares of the spin-one rotation matrix elements are
independent of whether the transverse rotation axis is chosen as \(x\) or
\(y\).  For an initial \(\sigma=0\),
\[
 P(0\to0)=\cos^2\Omega,\qquad
 P(0\to+1)=P(0\to-1)=\frac12\sin^2\Omega.
\]
For an initial \(\sigma=+1\),
\[
\begin{aligned}
 P(+1\to+1)&=\cos^4\frac{\Omega}{2},\\
 P(+1\to0)&=\frac12\sin^2\Omega,\\
 P(+1\to-1)&=\sin^4\frac{\Omega}{2}.
\end{aligned}
\]
Each set sums to one.  The phases of the amplitudes depend on the choice of
\(x\)- versus \(y\)-axis and on active/passive conventions, but these
probabilities do not.

\SupplementarySolution{21}{The Double Cover \(SL(2,\mathbb C)\)}
The determinant is
\[
 \det X=(x^0)^2-\mathbf x^2=-x_\mu x^\mu.
\]
For \(X'=AXA^\dagger\),
\[
 \det X'=\det A\,\det X\,\det A^\dagger=\det X,
\]
so the induced real linear map preserves the Minkowski norm.  Hermitian
\(X\) remains Hermitian, so real four-vectors map to real four-vectors.
Continuity from \(A=\mathbf1\) places the image in the proper
orthochronous component.

If \(AXA^\dagger=X\) for every Hermitian \(X\), then \(A\) commutes with
the full matrix algebra and is proportional to \(\mathbf1\).  The
determinant-one condition and the action on \(X=\mathbf1\) leave
\(A=\pm\mathbf1\).  Thus
\[
 SL(2,\mathbb C)/\{\pm\mathbf1\}\simeq SO^+(3,1).
\]

\SupplementarySolution{22}{Two Commuting Angular-Momentum Algebras}
Using
\([J_i,J_j]=i\epsilon_{ijk}J_k\),
\([J_i,K_j]=i\epsilon_{ijk}K_k\), and
\([K_i,K_j]=-i\epsilon_{ijk}J_k\), one finds
\[
 [A_i,A_j]=i\epsilon_{ijk}A_k,\qquad
 [B_i,B_j]=i\epsilon_{ijk}B_k,\qquad
 [A_i,B_j]=0.
\]
After complexification, finite-dimensional irreducible representations are
therefore tensor products labeled by
\((j_A,j_B)\), of dimension \((2j_A+1)(2j_B+1)\).

The Lorentz group is noncompact.  A faithful nontrivial finite-dimensional
unitary representation would have compact image in a unitary group, which
cannot represent the unbounded boost subgroup faithfully.  Correspondingly,
\(\mathbf A\) and \(\mathbf B\) are not separately Hermitian in these
finite multiplets.  Unitary particle representations instead act on
infinite-dimensional momentum-space Hilbert spaces.

\SupplementarySolution{23}{A Four-Vector as a Bispinor}
Write \(V=V^\mu\sigma_\mu\).  Under
\(A\in SL(2,\mathbb C)\),
\[
 V_{\alpha\dot\alpha}\longmapsto
 A_\alpha{}^\beta
 (A^\dagger)^{\dot\beta}{}_{\dot\alpha}
 V_{\beta\dot\beta},
\]
which is the tensor product of the undotted
\((\frac12,0)\) and dotted \((0,\frac12)\) representations:
\((\frac12,\frac12)\).  Expanding the resulting matrix uniquely in the
basis \(\{\sigma_\mu\}\) defines \(V'^\mu=\Lambda^\mu{}_\nu V^\nu\), and
S.2.21 proves that \(\Lambda\) is Lorentz.

A general bispinor has \(2\times2=4\) complex components, matching a
complex four-vector.  Imposing Hermiticity on \(V\) reduces these to four
real components and gives an ordinary real Lorentz vector.

\SupplementarySolution{24}{The Sign of a \(2\pi\) Rotation}
At \(\theta=2\pi\),
\[
 A(2\pi)=e^{-i\pi\sigma_3}=-\mathbf1,
\]
whereas at \(4\pi\), \(A(4\pi)=+\mathbf1\).  Both
\(\pm\mathbf1\) induce the same Lorentz transformation, so the image path
at \(2\pi\) is already the identity rotation although the lifted path in
\(SL(2,\mathbb C)\) has not closed.

In the spin-\(j\) representation,
\[
 U(R(2\pi))=(-1)^{2j}\mathbf1,\qquad U(R(4\pi))=\mathbf1.
\]
Thus half-integer spin gives a representation of the Lorentz group only up
to a sign, or equivalently an ordinary representation of its double cover.
This is the topological source of the familiar projective representation.

\SupplementarySolution{25}{The Galilean Mass Cocycle}
Let \(T(\mathbf a)\psi(\mathbf x,t)=\psi(\mathbf x-\mathbf a,t)\).
Using the stated boost law in the two possible orders gives
\[
 U(\mathbf v)T(\mathbf a)
 =e^{im\mathbf v\cdot\mathbf a}
   T(\mathbf a)U(\mathbf v).
\]
The extra phase cannot be seen in the spacetime action of the Galilei
group; it is a projective multiplier.  Expanding both sides to first order
in \(\mathbf v\) and \(\mathbf a\), with the generator conventions of
Chapter 2, yields
\[
 [K_i,P_j]=i\delta_{ij}M,\qquad M\ket{\psi}=m\ket{\psi}.
\]
The mass is central because the phase is a scalar and commutes with every
spacetime transformation.

\SupplementarySolution{26}{The Bargmann Group Commutator Loop}
The relation in S.2.25 immediately gives
\[
 U(\mathbf v)T(\mathbf a)U(\mathbf v)^{-1}T(\mathbf a)^{-1}
 =e^{im\mathbf v\cdot\mathbf a}\mathbf1.
\]
The associated sequence of coordinate transformations has zero net boost
and zero net translation, so it is the identity in the unextended Galilei
group.  Quantum mechanically it is a nontrivial element of the center.
Infinitesimally, the exponent is just the central
boost--translation commutator.

\SupplementarySolution{27}{Mass Superselection and Its Central-Extension
Escape}
For
\(\ket{\Psi}=c_1\ket{\psi_{m_1}}+c_2\ket{\psi_{m_2}}\), the loop gives
\[
 \ket{\Psi}\longmapsto
 c_1e^{im_1\mathbf v\cdot\mathbf a}\ket{\psi_{m_1}}
 +c_2e^{im_2\mathbf v\cdot\mathbf a}\ket{\psi_{m_2}}.
\]
Unless \(m_1=m_2\), this changes the relative phase even though the
underlying spacetime transformation is the identity.  Requiring the
identity to preserve every ray therefore forbids coherent superpositions
of different masses: Bargmann's mass-superselection rule.

In the centrally extended Galilei group, however, the loop is not declared
to be the identity; it is \(e^{i\alpha M}\).  If \(M\) is treated as an
operator with an additional conjugate degree of freedom, different mass
sectors are ordinary eigenspaces of the center.  The algebra alone then
does not forbid their superposition; a superselection rule requires an
additional physical restriction on observables or accessibility of that
central degree of freedom.

\SupplementarySolution{28}{The Galilean Algebra as a Poincar\'e Contraction}
Let \(H_{\rm rel}=cP^0=Mc^2+H_{\rm NR}+O(c^{-2})\) and
\(K_i^{\rm G}=J^{i0}/c\).  Restoring \(c\) in the Poincar\'e brackets gives
\[
\begin{aligned}
 [K_i^{\rm G},P_j]
 &=i\delta_{ij}\frac{H_{\rm rel}}{c^2}
 =i\delta_{ij}\left(M+\frac{H_{\rm NR}}{c^2}+\cdots\right),\\
 [K_i^{\rm G},H_{\rm NR}]&=iP_i+O(c^{-2}),\\
 [K_i^{\rm G},K_j^{\rm G}]
 &=-\frac{i}{c^2}\epsilon_{ijk}J_k.
\end{aligned}
\]
Taking \(c\to\infty\) leaves
\[
 [K_i^{\rm G},P_j]=i\delta_{ij}M,\qquad
 [K_i^{\rm G},H_{\rm NR}]=iP_i,\qquad
 [K_i^{\rm G},K_j^{\rm G}]=0.
\]
The relativistic rest energy becomes the central mass charge of the
contracted Galilean algebra.

\SupplementarySolution{29}{Massive Spin in \(2+1\) Dimensions}
Fixing \(k=(m,0,0)\) leaves only rotations of the two spatial coordinates,
so the connected little group is \(SO(2)\).  Its unitary characters are
\[
 D_s(\theta)=e^{is\theta}.
\]
For \(SO(2)\), \(D_s(\theta+2\pi)=D_s(\theta)\) requires
\(s\in\mathbb Z\).  On the double cover, a \(4\pi\) period permits
\(s\in\frac12\mathbb Z\).  The universal cover has
\(\theta\in\mathbb R\), so every \(s\in\mathbb R\) is allowed.

A spatial reflection conjugates \(R(\theta)\) to \(R(-\theta)\), and hence
\[
 D_s\longmapsto D_{-s}.
\]
Therefore a reflection-invariant theory must pair a nonzero \(s\) with
\(-s\), while a representation of the connected universal cover alone may
contain a single arbitrary real spin.

\SupplementarySolution{30}{Massless Little-Group Characters in \(2+1\)
Dimensions}
The stabilizer of \(k=(\kappa,0,\kappa)\) has one connected generator: a
null rotation \(N\), a particular linear combination of the single spatial
rotation and a transverse boost.  Its exponential satisfies
\[
 S(\alpha)S(\beta)=S(\alpha+\beta),
\]
so the connected little group is \((\mathbb R,+)\).
Every unitary irreducible representation is one-dimensional,
\[
 D_\rho(S(\alpha))=e^{i\rho\alpha},\qquad \rho\in\mathbb R.
\]
A reflection conjugates \(S(\alpha)\) to \(S(-\alpha)\), so it maps
\(\rho\) to \(-\rho\); nonzero \(\rho\) must be paired in a
reflection-invariant spectrum.  The representation used for ordinary
finite-polarization massless particles is the trivial character
\(\rho=0\).  Since the transverse rotation group in \(2+1\) dimensions is
trivial, this gives a single state at each null momentum rather than a
helicity multiplet.
\end{enumerate}
